\section{Runic Warriors in Faction Politics}

A Runic Warrior cannot hide their covenant. The engraved gorget, the carved staff, the tattoos curling up the neck—these are not easily concealed. This visibility shapes how factions perceive, use, and oppose them.

This chapter examines how major factions of Fate's Edge interact with Runic Warriors. Each entry includes the faction's general stance, specific points of tension or alliance, and sample story hooks.

\subsection{Chain-Lanterns of Thepyrgos}

The Chain-Lanterns hunt unlicensed magic. They prefer clear evidence: a Codex that can be seized, a Thiasos that can be caged, a confessor who can be made to speak.

A Runic Warrior confounds this preference.

\subsubsection{Stance}

Officially: Suspicious. A Runekeeper whose Codex is a tattoo cannot be disarmed. A Runekeeper whose Thiasos is a sentient blade cannot be separated from their familiar. This is a threat.

Unofficially: Pragmatic. A Runic Warrior who serves an approved Patron (Mykkiel, the Inquisitor Prime, the Oath of Flame \& Light) may be tolerated, even employed.

\subsubsection{Tensions}

\begin{itemize}
\item \textbf{Evidence.} How do you confiscate a tattoo? The Chain-Lanterns have resorted to branding over the tattoos, but this destroys the Rites and creates martyrs.
\item \textbf{Visibility.} A Runic Warrior cannot deny their covenant. This makes them easier to track—but also harder to prosecute, because their allegiance is public.
\item \textbf{The Bloodsworn.} Daughters of the Cord who serve the High Mother are technically licensed (the Cord has standing with the Thepyrgos courts). But a Bloodsworn who fights is something else. The Chain-Lanterns watch them closely.
\end{itemize}

\subsubsection{Alliances}

\begin{itemize}
\item \textbf{Inquisitor Prime.} Witch-hunters who serve the Prime are natural allies of the Chain-Lanterns. Their Codex is a censer. Their Thiasos is a lantern. They hunt together.
\item \textbf{Oath of Flame \& Light.} Dawn Warriors who swear vows before witnesses are legible to the Chain-Lanterns. A Dawn Knight who fights undead is useful. A Dawn Knight who questions procedure is not.
\end{itemize}

\subsubsection{Hooks}

\begin{itemize}
\item A Chain-Lantern prefect requests the party's help in identifying a rogue Bloodsworn. The suspect's Rites are tattooed, making them impossible to confiscate. The prefect wants the party to capture—not kill—so the tattoos can be studied.
\item A Runic Warrior is accused of unlicensed magic. Their Codex is a family heirloom, a shield passed down for generations. The Chain-Lanterns want to confiscate it. The Runekeeper wants the party to testify that the shield is not a Codex, but armor.
\item A Dawn Warrior has gone rogue. They still wear their engraved gorget. The Chain-Lanterns cannot remove it—the gorget is bolted to their breastplate. The party must decide: help the prefect capture the rogue, or help the rogue escape.
\end{itemize}

\subsection{Daughters of the Unbroken Cord}

The Cord deals in debts, not steel. But some Daughters have learned to fight. They are the Bloodsworn—Runic Warriors who serve the High Mother, their Rites tattooed into their flesh, their Thiasos a hollowed child or heartless bird.

\subsubsection{Stance}

Officially: The Bloodsworn are still Daughters. They pay their debts like any other sister.

Unofficially: Suspicion. A Daughter who fights may be less willing to collect gently. The Cord prefers subtlety. A Bloodsworn is not subtle.

\subsubsection{Tensions}

\begin{itemize}
\item \textbf{The Flesh Codex.} A Bloodsworn's Rites are tattooed on their skin. This is efficient—no Codex to lose—but it also means the Cord cannot audit their Rites without auditing their body.
\item \textbf{The Hollowed Thiasos.} A hollowed child or heartless bird is a witness, not a combatant. Bloodsworn who use their Thiasos in battle risk attracting the wrong kind of attention.
\item \textbf{The Open Palm, Closed Cord.} Some Bloodsworn develop the ability to store Fatigue as visible knots in their flesh. This is a gift from the High Mother. It is also a leash.
\end{itemize}

\subsubsection{Alliances}

\begin{itemize}
\item \textbf{The High Mother.} All Bloodsworn serve the High Mother, even if they also serve other Patrons. The Cord is their first debt.
\item \textbf{The Pale Shepherd.} Both the Cord and the Shepherd tend thresholds—birth, death, passage. Bloodsworn and Death Monks sometimes work together to escort the dying.
\end{itemize}

\subsubsection{Hooks}

\begin{itemize}
\item A Bloodsworn has stopped paying her tithe to the Cord. Her hollowed bird is silent. The Daughters want the party to find her—and remind her of her debt.
\item A rogue Bloodsworn is using her Flesh Codex to record forbidden Rites. The Cord cannot audit her because her tattoos are her own. The party must capture her so a Weaver can read her skin.
\item A Bloodsworn offers to help the party with a dangerous mission. Her price: a witness to her next payment to the Cord. She wants someone to know she is still loyal.
\end{itemize}

\subsection{Aeler Vent-Priors and Correctors}

The Aeler measure everything. A Runic Warrior who is also efficient is respected. A Runic Warrior who is inefficient is a liability.

\subsubsection{Stance}

Officially: Runic Warriors are evaluated on their utility. An Artificer who maintains the Underways is valuable. A Cairn-Knight who guards a burial mound is acceptable. A Beast Warrior who cannot keep ledgers is not.

Unofficially: Fascination. The Aeler appreciate a Codex that is also a tool. A shield that records Rites is efficient. A staff that grows its own runes is elegant.

\subsubsection{Tensions}

\begin{itemize}
\item \textbf{The Breath Bond.} Aeler Correctors use breath-bonds to secure contracts. A Runic Warrior whose Codex is part of their body cannot be breath-bonded in the usual way. This confuses the Correctors.
\item \textbf{The Forge Clause.} Clockwork Monad Artificers sometimes engrave their Rites onto their own prosthetics. This is allowed—but the Correctors audit the prosthetics every season.
\item \textbf{The Cairn Oath.} Cairn-Knights who guard Aeler burial mounds are respected, but they are also watched. The ancestors speak through them. The Vent-Priors want to know what the ancestors are saying.
\end{itemize}

\subsubsection{Alliances}

\begin{itemize}
\item \textbf{Clockwork Monad.} Artificers who serve the Monad are natural partners of the Vent-Priors. Efficiency is efficiency.
\item \textbf{Varnek Karn.} Cairn-Knights who serve Varnek Karn sometimes negotiate with Aeler ancestors. The Vent-Priors tolerate this because it prevents hauntings.
\end{itemize}

\subsubsection{Hooks}

\begin{itemize}
\item An Aeler Vent-Prior hires the party to audit a Clockwork Monad Artificer. The Artificer's Codex is a prosthetic arm. The Vent-Prior wants to know if the arm has been modified without authorization.
\item A Cairn-Knight refuses to allow Aeler surveyors near her burial mound. She claims the ancestors are restless. The Vent-Priors want the party to negotiate access—or confirm that the ancestors are real.
\item An Artificer's sentient Thiasos (a clockwork construct) has begun speaking in a language no one understands. The Vent-Priors suspect unauthorized modifications. The Artificer insists the construct is evolving on its own.
\end{itemize}

\subsection{Lethai Courts (Crimson Valley and City Elves)}

The Lethai do not send armies. They send Ghe'hai—elite operatives who are warriors, diplomats, and assassins all at once. A Ghe'hai operative is a Runic Warrior, their Codex a woven mask, their Thiasos a spider-spirit.

\subsubsection{Stance}

Officially: The Ghe'hai are essential. They protect Lethai interests without drawing attention. A Ghe'hai who fails is replaced. A Ghe'hai who succeeds is never named.

Unofficially: Fear. The Ghe'hai are weapons. Weapons can be turned. The courts watch them closely.

\subsubsection{Tensions}

\begin{itemize}
\item \textbf{The Mask.} A Ghe'hai's Codex is a mask. When worn, they are the operative. When removed, they are no one. This ambiguity is useful—and unsettling.
\item \textbf{The Spider-Spirit.} A Ghe'hai's Thiasos weaves invisible threads. These threads can bind, trip, or kill. The courts appreciate the utility. They also worry about whose neck the thread is wrapped around.
\item \textbf{The Unspoken Question.} Some Ghe'hai develop the ability to force a truthful answer from a target. This is a Rite of Inaea. The courts approve—when it is used on enemies.
\end{itemize}

\subsubsection{Alliances}

\begin{itemize}
\item \textbf{Inaea.} Most Ghe'hai serve Inaea, the Spider's Mercy. Their threads bind gently—or not.
\item \textbf{Isoka.} Some Ghe'hai serve Isoka, the Serpent's Shedding. They change identities as easily as changing masks.
\item \textbf{Ikasha.} A few serve Ikasha, the Shadow's Silence. They are the ones no one sees coming.
\end{itemize}

\subsubsection{Hooks}

\begin{itemize}
\item A Ghe'hai operative has gone silent. The Lethai courts want the party to find them—and determine if they have been turned.
\item A rival court has deployed its own Ghe'hai. The party is caught between two operatives, each wearing a different mask, each claiming to represent the true Lethai interest.
\item A Ghe'hai operative offers to help the party with a mission. Their price: a witness to their next assignment. They want someone to know what they did—in case they do not return.
\end{itemize}

\subsection{Black Banner Companies}

Mercenary companies care about results. A Runic Warrior who can fight is valuable. A Runic Warrior who cannot be disarmed is invaluable.

\subsubsection{Stance}

Officially: Runic Warriors are hired like any other soldier. Their Codex is their equipment. Their Thiasos is their mount or weapon.

Unofficially: Respect. A warrior who carries their covenant on their body has proven something. The Black Banners do not ask what. They only ask if the warrior can hold the line.

\subsubsection{Tensions}

\begin{itemize}
\item \textbf{Payday.} A Runic Warrior whose Codex is a tattoo cannot sell their equipment. This makes them less likely to desert—but also harder to recruit, because they have no collateral.
\item \textbf{The Condotta Flip.} A mercenary captain who switches sides may find that their Runic Warrior refuses to follow. Oaths carved into flesh are harder to break than oaths written on paper.
\item \textbf{The Blood Price.} A Cairn-Knight who serves Varnek Karn may demand wergild for fallen comrades. This is unusual in mercenary companies. The captain must decide: pay, or lose the knight.
\end{itemize}

\subsubsection{Alliances}

\begin{itemize}
\item \textbf{Varnek Karn.} Cairn-Knights who serve Varnek Karn sometimes hire on as mercenaries. They are expensive—they demand wergild—but they are loyal.
\item \textbf{The Savage Heart.} Beast Warriors sometimes join mercenary companies for the hunt. They leave when the hunt ends.
\end{itemize}

\subsubsection{Hooks}

\begin{itemize}
\item A mercenary captain hires the party to recruit a Runic Warrior. The target is a Cairn-Knight who refuses to leave their burial mound. The party must convince them—or find another.
\item A Runic Warrior in a mercenary company has been accused of desertion. They claim their Codex (a sentient blade) ordered them to leave. The party must determine who is lying.
\item A mercenary company is being hunted by a Cursed Knight (Malachai). The knight was once a member. Now they cannot stop killing. The company wants the party to stop them—permanently.
\end{itemize}

\subsection{The Covenant (Mykkielan Courts)}

The Covenant deals in law, not war. But some Runekeepers serve Mykkiel as Justicars—warriors who enforce oaths with steel.

\subsubsection{Stance}

Officially: Justicars are officers of the court. Their Codex is a shield. Their Thiasos is a spirit-hound that tracks oath-breakers. They are respected—and feared.

Unofficially: Distrusted. A Justicar who enforces law with a sword is a reminder that the Covenant's power rests on violence.

\subsubsection{Tensions}

\begin{itemize}
\item \textbf{The Shield Codex.} A Justicar's shield records their Rites. This is practical—the shield is always in hand—but it also means the Justicar cannot be disarmed without losing their Rites.
\item \textbf{The Spirit-Hound.} A Justicar's Thiasos tracks oath-breakers. This is useful in court. It is less useful when the hound decides that a judge has broken an oath.
\item \textbf{The Sentence.} Some Justicars develop the ability to speak a binding judgment. The target is compelled to obey—or suffer. The Covenant approves, as long as the judgment is lawful.
\end{itemize}

\subsubsection{Alliances}

\begin{itemize}
\item \textbf{Mykkiel.} All Justicars serve Mykkiel, the Arbiter of the Covenant.
\item \textbf{The Witness.} Justicars sometimes work with Truth-Sworn to verify testimony.
\end{itemize}

\subsubsection{Hooks}

\begin{itemize}
\item A Justicar has been accused of enforcing an unlawful oath. The Covenant wants the party to investigate. The Justicar insists the oath was binding.
\item A spirit-hound has gone rogue. It is tracking oath-breakers who have already paid their debts. The Justicar cannot control it. The party must help.
\item A Justicar offers to help the party with a legal dispute. Their price: a witness to their next judgment. They want someone to verify that they acted lawfully.
\end{itemize}

\subsection{The Unaffiliated}

Many Runic Warriors belong to no faction. They are ronin, wanderers, hermits. They owe allegiance to their Patron—and to no one else.

\subsubsection{Stance}

Officially: Suspicious. A Runekeeper with no faction is a Runekeeper with no accountability.

Unofficially: Useful. An unaffiliated Runic Warrior can be hired for deniable work.

\subsubsection{Tensions}

\begin{itemize}
\item \textbf{No Witness.} Unaffiliated Runekeepers often lack a witness for their Rites. This makes their magic less stable—and more dangerous.
\item \textbf{No Backing.} An unaffiliated Runekeeper has no faction to protect them. If they are accused of a crime, they are on their own.
\end{itemize}

\subsubsection{Alliances}

None. That is the point.

\subsubsection{Hooks}

\begin{itemize}
\item An unaffiliated Runic Warrior approaches the party with a job. They need witnesses for a High Rite. They will pay in service, not coin.
\item A rogue Bloodsworn has gone unaffiliated. The Cord wants her back. She wants the Cord to forget her. The party is caught in the middle.
\item An unaffiliated Beast Warrior has been hunting in territory claimed by a mercenary company. The company wants the party to drive the warrior off. The warrior wants the party to witness their next hunt.
\end{itemize}

The thread held. The water carried. That is not a Rite. That is grace.
